\documentclass[tikz,border=8pt]{standalone}
\usepackage{tikz}
\usepackage{amsmath}
\usepackage{amssymb}
\usepackage{ctex}
\usetikzlibrary{calc, positioning, arrows.meta, decorations.pathreplacing}

\begin{document}
\begin{tikzpicture}[>=Stealth]

%% ===== 左图：标准基 e1, e2 =====
\begin{scope}[xshift=0cm]

  % 标准基网格
  \draw[gray!25, thin] (-0.5,-0.5) grid (4.5,4.5);
  \draw[->, gray!60, thin] (-0.5,0) -- (4.8,0) node[right, font=\scriptsize] {$x$};
  \draw[->, gray!60, thin] (0,-0.5) -- (0,4.8) node[above, font=\scriptsize] {$y$};

  % 标准基网格线（每步 1）
  \foreach \i in {0,1,2,3,4}{
    \draw[blue!15, thin] (\i,-0.5) -- (\i,4.5);
    \draw[blue!15, thin] (-0.5,\i) -- (4.5,\i);
  }

  % 标准基向量
  \draw[->, very thick, blue!80!black] (0,0) -- (1,0)
    node[below, font=\small] {$\mathbf{e}_1=(1,0)$};
  \draw[->, very thick, red!70!black] (0,0) -- (0,1)
    node[left, font=\small] {$\mathbf{e}_2=(0,1)$};

  % 示例点：用标准基坐标表示
  \fill[gray!70] (3,2) circle (2pt)
    node[above right, font=\scriptsize] {$3\mathbf{e}_1+2\mathbf{e}_2=(3,2)$};
  \draw[->, gray!50, dashed, thin] (0,0) -- (3,2);

  % 原点
  \fill (0,0) circle (1.8pt) node[below left, font=\scriptsize] {$O$};

  % 子图标题
  \node[font=\small\bfseries] at (2.0, -0.9) {标准基 $\{\mathbf{e}_1,\mathbf{e}_2\}$};
  \node[font=\scriptsize, fill=white, draw=gray!50, thin, rounded corners=2pt, inner sep=3pt]
    at (2.0, -1.55) {正交单位基，坐标即分量};

\end{scope}

%% ===== 右图：一般基 u1=(1,1), u2=(1,-1) =====
\begin{scope}[xshift=6.5cm]

  % 背景网格（标准网格参考）
  \draw[gray!15, thin] (-0.5,-3.5) grid (4.5,4.5);
  \draw[->, gray!60, thin] (-0.5,0) -- (4.8,0) node[right, font=\scriptsize] {$x$};
  \draw[->, gray!60, thin] (0,-3.5) -- (0,4.8) node[above, font=\scriptsize] {$y$};

  % 一般基斜坐标网格（u1=(1,1), u2=(1,-1) 方向）
  \foreach \i in {0,1,2,3}{
    % 沿 u1=(1,1) 方向的网格线（过 i*u2 点）
    \draw[teal!20, thin] ({\i*1-0.5*1}, {\i*(-1)-0.5*(-1)}) -- ({\i*1+3.5*1}, {\i*(-1)+3.5*(-1)});
    % 沿 u2=(1,-1) 方向的网格线（过 i*u1 点）
    \draw[teal!20, thin] ({\i*1-1.5*1}, {\i*1-1.5*(-1)}) -- ({\i*1+3*1}, {\i*1+3*(-1)});
  }
  % 补充负方向
  \foreach \i in {1,2}{
    \draw[teal!20, thin] ({-\i*1-0.5}, {-\i*(-1)+0.5}) -- ({-\i*1+3.5}, {-\i*(-1)-3.5});
  }

  % 一般基向量
  \draw[->, very thick, teal!80!black] (0,0) -- (1,1)
    node[above left, font=\small] {$\mathbf{u}_1=(1,1)$};
  \draw[->, very thick, orange!80!black] (0,0) -- (1,-1)
    node[below left, font=\small] {$\mathbf{u}_2=(1,-1)$};

  % 示例点：用 u1,u2 坐标
  \fill[gray!70] (3,1) circle (2pt)
    node[above right, font=\scriptsize] {$2\mathbf{u}_1+\mathbf{u}_2=(3,1)$};
  \draw[->, gray!50, dashed, thin] (0,0) -- (3,1);

  % 原点
  \fill (0,0) circle (1.8pt) node[below left, font=\scriptsize] {$O$};

  % 子图标题
  \node[font=\small\bfseries] at (2.0, -1.7) {一般基 $\{\mathbf{u}_1,\mathbf{u}_2\}$};
  \node[font=\scriptsize, fill=white, draw=gray!50, thin, rounded corners=2pt, inner sep=3pt]
    at (2.0, -2.35) {斜坐标基，仍张成 $\mathbb{R}^2$};

\end{scope}

%% 中间标注：维数不变
\node[font=\normalsize\bfseries, fill=white, draw=gray!60, thin, rounded corners=4pt,
      inner sep=6pt] at (5.25, 5.8)
  {无论选哪组基，维数 $= 2$（$\mathbb{R}^2$ 是 2 维空间）};

%% 整体标题
\node[font=\large\bfseries] at (5.25, 6.8)
  {基与维数：同一平面上的不同基};

\end{tikzpicture}
\end{document}
